% Options for packages loaded elsewhere
\PassOptionsToPackage{unicode}{hyperref}
\PassOptionsToPackage{hyphens}{url}
\PassOptionsToPackage{dvipsnames,svgnames,x11names}{xcolor}
\documentclass[
  11pt,
  twocolumn]{article}
\usepackage{xcolor}
\usepackage[margin=1.5cm]{geometry}
\usepackage{amsmath,amssymb}
\setcounter{secnumdepth}{-\maxdimen} % remove section numbering
\usepackage{iftex}
\ifPDFTeX
  \usepackage[T1]{fontenc}
  \usepackage[utf8]{inputenc}
  \usepackage{textcomp} % provide euro and other symbols
\else % if luatex or xetex
  \usepackage{unicode-math} % this also loads fontspec
  \defaultfontfeatures{Scale=MatchLowercase}
  \defaultfontfeatures[\rmfamily]{Ligatures=TeX,Scale=1}
\fi
\usepackage{lmodern}
\ifPDFTeX\else
  % xetex/luatex font selection
  \setmainfont[]{Arial}
\fi
% Use upquote if available, for straight quotes in verbatim environments
\IfFileExists{upquote.sty}{\usepackage{upquote}}{}
\IfFileExists{microtype.sty}{% use microtype if available
  \usepackage[]{microtype}
  \UseMicrotypeSet[protrusion]{basicmath} % disable protrusion for tt fonts
}{}
\makeatletter
\@ifundefined{KOMAClassName}{% if non-KOMA class
  \IfFileExists{parskip.sty}{%
    \usepackage{parskip}
  }{% else
    \setlength{\parindent}{0pt}
    \setlength{\parskip}{6pt plus 2pt minus 1pt}}
}{% if KOMA class
  \KOMAoptions{parskip=half}}
\makeatother
\usepackage{graphicx}
\makeatletter
\newsavebox\pandoc@box
\newcommand*\pandocbounded[1]{% scales image to fit in text height/width
  \sbox\pandoc@box{#1}%
  \Gscale@div\@tempa{\textheight}{\dimexpr\ht\pandoc@box+\dp\pandoc@box\relax}%
  \Gscale@div\@tempb{\linewidth}{\wd\pandoc@box}%
  \ifdim\@tempb\p@<\@tempa\p@\let\@tempa\@tempb\fi% select the smaller of both
  \ifdim\@tempa\p@<\p@\scalebox{\@tempa}{\usebox\pandoc@box}%
  \else\usebox{\pandoc@box}%
  \fi%
}
% Set default figure placement to htbp
\def\fps@figure{htbp}
\makeatother
\setlength{\emergencystretch}{3em} % prevent overfull lines
\providecommand{\tightlist}{%
  \setlength{\itemsep}{0pt}\setlength{\parskip}{0pt}}
\pagenumbering{gobble}
\setlength{\columnsep}{2em}
\setlength{\headheight}{13.6pt}
\addtolength{\topmargin}{-15pt}
\usepackage{fancyhdr}
\usepackage{hyperref}
\pagestyle{fancy}
\fancyfoot[L]{\normalsize \color{red}▲ \color{black}Red triangles indicate parameters that exceeded the acceptable threshold at least once in 2024.\\An \textbf{accessible HTML version} of this document is available in the \textbf{\href{https://irma.nps.gov/DataStore/Reference/Profile/2316392}{{NPS DataStore}}}.}
\fancypagestyle{plain}{\pagestyle{fancy}}
\renewcommand{\headrulewidth}{0pt}
\usepackage{bookmark}
\IfFileExists{xurl.sty}{\usepackage{xurl}}{} % add URL line breaks if available
\urlstyle{same}
\hypersetup{
  colorlinks=true,
  linkcolor={Maroon},
  filecolor={Maroon},
  citecolor={Blue},
  urlcolor={RoyalBlue},
  pdfcreator={LaTeX via pandoc}}

\title{\vspace{-2em} \huge \textbf{2025 Water Quality Summary}\\
\LARGE Antietam National Battlefield: Sharpsburg Creek \vspace{-1em}}
\author{}
\date{\vspace{-2.5em}}

\begin{document}
\maketitle

\Large \textbf{Percent of Measurements Passing} \vspace{-1.5em}

\begin{flushleft}\includegraphics[width=1.2\linewidth]{pdfs_2025/ANTI_SHCK_2025_summary_files/figure-latex/print figure-1} \end{flushleft}

\Large \textbf{Acid Neutralizing Capacity (ANC)}
\textbf{\textcolor{red}{ }}\\
\normalsize \textbf{Acceptable values:} Above 600 µeq/L\\
\textbf{2025 range:} 4822 -- 5322.15 µeq/L\\
All ANC measurements at Sharpsburg Creek were passing in 2025 and the
five years prior.\\
\vspace{1pt}\\
\Large \textbf{Dissolved Oxygen (DO)} \textbf{\textcolor{red}{ }}\\
\normalsize \textbf{Acceptable values:} Above 5 mg/L\\
\textbf{2025 range:} 9.55 -- 11.33 mg/L\\
All DO measurements at Sharpsburg Creek were passing in 2025 and the
five years prior.\\
\vspace{1pt}\\
\Large \textbf{pH} \textbf{\textcolor{red}{ }}\\
\normalsize \textbf{Acceptable range:} 6 -- 9\\
\textbf{2025 range:} 8.11 -- 8.59\\
All pH measurements at Sharpsburg Creek were passing in 2025 and the
five years prior.\\
\vspace{1pt}\\
\Large \textbf{Specific Conductance} \textbf{\textcolor{red}{▲}}\\
\normalsize \textbf{Acceptable values:} Below 171 µS/cm\\
\textbf{2025 range:} 661 -- 730 µS/cm\\
All specific conductance measurements at Sharpsburg Creek exceeded the
acceptable threshold in 2025 and the five years prior.

\vfill\eject

\[\\[-1.55em]\]

\begin{flushright}\includegraphics[width=0.85\linewidth,alt={test}]{Site images/ANTI_SHCK} \end{flushright}

\vspace{-1.5em}
\begin{flushright}
\begin{minipage}{0.85\columnwidth}
\small Sharpsburg Creek is a warmwater karst stream.
\end{minipage}
\end{flushright}

\[\\[-2.85em]\]

\Large \textbf{Total Nitrogen} \textbf{\textcolor{red}{▲}}\\
\normalsize \textbf{Acceptable values:} Below 0.31 mg/L\\
\textbf{2025 range:} 4.72 -- 5.8 mg/L\\
All total nitrogen measurements at Sharpsburg Creek exceeded the
acceptable threshold in 2025 and the five years prior.\\
\vspace{-1pt}\\
\Large \textbf{Total Phosphorus} \textbf{\textcolor{red}{▲}}\\
\normalsize \textbf{Acceptable values:} Below 0.01 mg/L\\
\textbf{2025 range:} 0.022 -- 0.031 mg/L\\
All total phosphorus measurements at Sharpsburg Creek exceeded the
acceptable threshold in 2025 and the five years prior.\\
\vspace{-1pt}\\
\Large \textbf{Water Temperature} \textbf{\textcolor{red}{ }}\\
\normalsize \textbf{Acceptable values:} Below 32 °C\\
\textbf{2025 range:} 9.2 -- 16.3 °C\\
All water temperature measurements at Sharpsburg Creek were passing in
2025 and the five years prior.\\
\vspace{2pt}\\
\Large \textbf{Water Quality Monitoring in the NCR}\\
\normalsize This report summarizes water quality data from 2020 to 2025
for seven parameters. The National Capital Region Inventory \&
Monitoring Network (NCRN I\&M) monitors additional water quality
parameters, including depth, discharge, conductivity, and total
dissolved solids. These data are available in the
\href{https://irma.nps.gov/DataStore/Reference/Profile/2237386}{NCRN
I\&M water quality data visualizer}. Visit the
\href{https://www.nps.gov/im/ncrn/water.htm}{NCRN I\&M water quality
monitoring webpage} to learn more.

\end{document}
